% ===========================================================
%  บทสรุปผู้บริหาร (2 หน้า) — กลุ่มแผนงานใต้ร่มพระบารมี 2569
%  ข้อมูล ณ 4 สิงหาคม 2569 · ตัวเลขทุกจุดมาจากฐานข้อมูลติดตามโครงการ
%  build: xelatex exec-brief.tex
% ===========================================================
\documentclass[11pt]{article}

\usepackage{fontspec}
\usepackage{polyglossia}
\setmainlanguage{thai}
\setotherlanguage{english}

\IfFontExistsTF{TH Sarabun New}{\setmainfont{TH Sarabun New}[Scale=1.0]}{%
  \IfFontExistsTF{TH SarabunPSK}{\setmainfont{TH SarabunPSK}[Scale=1.0]}{\setmainfont{Sarabun}[Scale=1.0]}}

% --- การตัดคำไทย (ค่าเดียวกับเล่มรายงาน) ---
\XeTeXlinebreaklocale "th"
\XeTeXlinebreakskip = 0pt plus 0.04em minus 0.02em
\XeTeXlinebreakpenalty = 120
\tolerance=700
\emergencystretch=3em
\hbadness=10000

\usepackage[a4paper,top=1.5cm,bottom=1.4cm,left=1.7cm,right=1.7cm]{geometry}
\usepackage[originalparindent]{ragged2e}
\usepackage{xcolor}
\usepackage{tikz}
\usepackage{booktabs}
\usepackage{colortbl}
\usepackage{array}
\usepackage{float}
\usepackage{enumitem}
\usepackage{fancyhdr}

\definecolor{FirstBlue}{HTML}{2D4E9C}
\definecolor{SecondBlue}{HTML}{4A96CD}
\definecolor{rpfgold}{HTML}{C79A3C}
\definecolor{headerblue}{RGB}{213, 226, 243}
\definecolor{rowblue}{RGB}{231, 238, 249}
\definecolor{tipgreen}{HTML}{2E7D32}
\definecolor{warnred}{HTML}{C0392B}
\definecolor{greybox}{HTML}{F4F6F9}

\newcolumntype{L}[1]{>{\RaggedRight\arraybackslash}p{#1}}
\newcolumntype{C}[1]{>{\centering\arraybackslash}p{#1}}

\setlength{\parindent}{0pt}
\setlength{\parskip}{4pt}
\renewcommand{\arraystretch}{1.25}

\pagestyle{fancy}\fancyhf{}
\renewcommand{\headrulewidth}{0pt}
\fancyfoot[L]{\footnotesize\color{black!55}กลุ่มแผนงานใต้ร่มพระบารมี มทร.ล้านนา $\cdot$ ข้อมูล ณ 4 สิงหาคม 2569}
\fancyfoot[R]{\footnotesize\color{black!55}หน้า \thepage/2}

% ---------- ตัวเลขเด่น ----------
\newcommand{\bignum}[3]{%
  \begin{tikzpicture}
    \node[draw=SecondBlue!35,line width=0.6pt,fill=#3,rounded corners=3pt,
          minimum width=3.55cm,minimum height=1.5cm,inner sep=3pt] {%
      \begin{minipage}{3.3cm}\centering
        {\color{FirstBlue}\fontsize{19}{21}\selectfont\bfseries #1}\par\vspace{1pt}
        {\fontsize{9.5}{11}\selectfont #2}
      \end{minipage}};
  \end{tikzpicture}}

% ---------- หัวข้อ ----------
\newcommand{\sechead}[1]{%
  \vspace{5pt}
  {\color{FirstBlue}\fontsize{13.5}{16}\selectfont\bfseries #1}%
  \par\vspace{-3pt}{\color{rpfgold}\rule{\linewidth}{1.2pt}}\par\vspace{2pt}}

% กล่องไฮไลต์ — แถบสีซ้าย + พื้นอ่อน (ไม่พึ่ง tcolorbox)
\newsavebox{\flagbuf}
\newenvironment{flagbox}[1]{%
  \colorlet{flagbar}{#1}%
  \par\vspace{3pt}\noindent
  \begin{lrbox}{\flagbuf}%
  \begin{minipage}{\dimexpr\linewidth-0.85cm\relax}\RaggedRight\small\vspace{3pt}}%
  {\vspace{3pt}\end{minipage}\end{lrbox}%
  \colorbox{flagbar!6}{\makebox[\dimexpr\linewidth-2\fboxsep\relax][l]{%
    \textcolor{flagbar}{\rule[-\dimexpr\dp\flagbuf+3pt\relax]{2.6pt}{\dimexpr\ht\flagbuf+\dp\flagbuf+6pt\relax}}%
    \hspace{5pt}\usebox{\flagbuf}}}\par\vspace{3pt}}

\begin{document}
\RaggedRight

% ======================= หน้า 1 =======================
\begin{tikzpicture}[remember picture,overlay]
  \fill[FirstBlue] (current page.north west) rectangle ([yshift=-3.9cm]current page.north east);
  \fill[rpfgold]   ([yshift=-3.9cm]current page.north west) rectangle ([yshift=-4.03cm]current page.north east);
\end{tikzpicture}
\vspace*{-0.15cm}
{\color{white}\fontsize{11}{13}\selectfont บทสรุปผู้บริหาร $\cdot$ รายงานความก้าวหน้า รอบ 10~เดือน}\par\vspace{3pt}
{\color{white}\fontsize{23}{26}\selectfont\bfseries กลุ่มแผนงานใต้ร่มพระบารมี}\par\vspace{2pt}
{\color{rpfgold}\fontsize{14}{16}\selectfont\bfseries มหาวิทยาลัยเทคโนโลยีราชมงคลล้านนา}\hfill
{\color{white!85}\fontsize{10.5}{12}\selectfont ปีงบประมาณ 2569 $\cdot$ 1 ต.ค. 2568 -- 4 ส.ค. 2569}
\vspace{0.85cm}

\begin{center}
\bignum{7,986,583}{งบประมาณจัดสรร (บาท)}{white}\hspace{-3pt}
\bignum{4,460,332}{เบิกจ่ายแล้ว $=$ 56\%}{rowblue}\hspace{-3pt}
\bignum{61}{โครงการย่อย ใน 3 ง8}{white}\hspace{-3pt}
\bignum{55}{วันเหลือถึงสิ้นปีงบ}{rowblue}
\end{center}
\vspace{-4pt}

\sechead{1. สถานะโดยสรุป}

แผนงานดำเนินการ 61~โครงการ ภายใต้โครงการหลัก 3 โครงการตามแบบ ง8
โดยหน่วยงานภายในมหาวิทยาลัย 11~หน่วยงาน หัวหน้าโครงการ 42~ท่าน
ครอบคลุมพื้นที่ดำเนินงาน 48~พื้นที่ ในศูนย์พัฒนาโครงการหลวง สถานีเกษตรหลวง
โครงการตามพระราชดำริ และชุมชนบนพื้นที่สูง

\begin{flagbox}{tipgreen}
\textbf{ทำได้เกินเป้าแล้ว} --- ตัวชี้วัดที่ 39 (องค์ความรู้ เทคโนโลยี และนวัตกรรมที่นำไปใช้
ในโครงการหลวงและโครงการตามพระราชดำริ) ซึ่งเป็น\textbf{เป้าหมายโดยตรงของแผนงาน}
รับมาดำเนินการ \textbf{62 จากเป้า 60 คิดเป็น 103\%} กระจายใน 37~โครงการ ;
อนุมัติโครงการแล้ว 59 จาก 61 ; โอนงบถึงผู้รับผิดชอบแล้ว 98\%
\end{flagbox}

\begin{flagbox}{warnred}
\textbf{ต้องเร่งด่วน} --- เบิกจ่าย \textbf{56\%} ขณะที่เวลาปีงบประมาณผ่านไป \textbf{85\%}
\textbf{ช้ากว่ากรอบเวลา 29~จุดร้อยละ} ต้องเบิกจ่ายเพิ่มอีก \textbf{3,526,251~บาท ใน 55~วัน}
โดยมีใบสั่งซื้อรองรับแล้วเพียง 825,156~บาท เท่ากับวงเงินกว่า \textbf{2.7 ล้านบาทยังไม่เริ่มกระบวนการจัดซื้อจัดจ้าง}
\end{flagbox}

\sechead{2. ตัวเลขสำคัญ}

\begin{minipage}[t]{0.485\textwidth}
{\footnotesize\textbf{\color{FirstBlue}สถานะงบประมาณภาพรวม}}\par\vspace{2pt}
\small
\begin{tabular}{L{3.55cm}rr}
\toprule
\rowcolor{headerblue}
\textbf{รายการ} & \textbf{บาท} & \textbf{\%}\\
\midrule
งบประมาณจัดสรร & 7,986,583 & 100\\
\rowcolor{rowblue}
โอนถึงผู้รับผิดชอบ & 7,798,583 & 98\\
ผูกพันใบสั่งซื้อ & 825,156 & 10\\
\rowcolor{rowblue}
\textbf{เบิกจ่ายจริง} & \textbf{4,460,332} & \textbf{56}\\
เบิกจ่าย$+$ผูกพัน & 5,285,488 & 66\\
\rowcolor{rowblue}
\textbf{คงเหลือต้องเบิก} & \textbf{3,526,251} & \textbf{44}\\
\bottomrule
\end{tabular}
\end{minipage}\hfill
\begin{minipage}[t]{0.485\textwidth}
{\footnotesize\textbf{\color{FirstBlue}รายโครงการหลัก ง8}}\par\vspace{2pt}
\small
\begin{tabular}{L{2.75cm}C{0.75cm}rr}
\toprule
\rowcolor{headerblue}
\textbf{โครงการหลัก} & \textbf{จำนวน} & \textbf{เบิกจ่าย} & \textbf{\%}\\
\midrule
ง8-1 เทคโนโลยี นวัตกรรมสู่ชุมชน & 16 & 1,064,116 & 54\\
\rowcolor{rowblue}
ง8-2 องค์ความรู้ ยกระดับคุณภาพชีวิต & 14 & 1,360,534 & 68\\
ง8-3 พัฒนากำลังคน ลดความเหลื่อมล้ำ & 31 & 2,035,682 & 51\\
\midrule
\rowcolor{headerblue}
\textbf{รวม} & \textbf{61} & \textbf{4,460,332} & \textbf{56}\\
\bottomrule
\end{tabular}
\end{minipage}

\vspace{4pt}
{\footnotesize\color{black!60}ร้อยละการเบิกจ่ายทุกจุดคำนวณจาก\textbf{งบประมาณที่ได้รับจัดสรร} เพื่อให้ตัวเลขรายโครงการหลักบวกกลับตรงกับภาพรวม}

\sechead{3. ผลตัวชี้วัดที่รับมาดำเนินการ}

{\small\renewcommand{\arraystretch}{1.08}
\begin{tabular}{L{9.4cm}C{2.3cm}C{1.5cm}L{2.4cm}}
\toprule
\rowcolor{headerblue}
\textbf{ตัวชี้วัด} & \textbf{รับมา/เป้า} & \textbf{ร้อยละ} & \textbf{สถานะ}\\
\midrule
ที่ 39 องค์ความรู้สู่โครงการหลวงและโครงการตามพระราชดำริ & 62 / 60 & 103\% & {\color{tipgreen}\bfseries ผ่านเป้า}\\
\rowcolor{rowblue}
ที่ 35 กำลังคนที่ได้รับการพัฒนาทักษะ & 288 / 400 & 72\% & ตามแผน\\
ที่ 40 องค์ความรู้ยกระดับคุณภาพชีวิตชุมชน & 45 / 100 & 45\% & ต้องเร่ง\\
\rowcolor{rowblue}
ที่ 10 เครือข่ายความร่วมมือ & --- & 18\% & {\color{warnred}\bfseries เชิงโครงสร้าง}\\
ที่ 17 ทรัพย์สินทางปัญญา & --- & 13\% & {\color{warnred}\bfseries เชิงโครงสร้าง}\\
\rowcolor{rowblue}
ที่ 38 ความร่วมมือกับสถานประกอบการ & --- & 8\% & {\color{warnred}\bfseries เชิงโครงสร้าง}\\
\bottomrule
\end{tabular}}

\vspace{3pt}
\begin{flagbox}{SecondBlue}
\textbf{ข้อควรทราบในการอ่านตัวเลข} --- ยังไม่มีการบันทึกรายงานผลรายกิจกรรมเข้าระบบติดตาม
(0 จาก 268~กิจกรรม) ตัวเลขตัวชี้วัดข้างต้นจึงเป็น\textbf{ค่าที่โครงการรับมาดำเนินการ}
มิใช่ผลสำเร็จที่ผ่านการตรวจรับ การปิดช่องว่างนี้เป็นเงื่อนไขจำเป็นของรายงานฉบับสมบูรณ์
\end{flagbox}

% ======================= หน้า 2 =======================
\sechead{4. ประเด็นวิเคราะห์}

\begin{enumerate}[leftmargin=1.5em,itemsep=2pt,topsep=2pt]
  \item \textbf{ปัญหาอยู่ที่กระบวนการ มิใช่เนื้องาน} --- ตัวชี้วัดภารกิจหลักทำได้เกินเป้าแล้ว
        แต่การเบิกจ่ายและการรายงานผลตามไม่ทัน มาตรการที่ควรใช้จึงเป็นด้านการบริหารจัดการ
        และการสนับสนุนกระบวนการพัสดุ มิใช่การปรับเปลี่ยนเนื้องานของโครงการ
  \item \textbf{คอขวดไม่ได้อยู่ต้นทาง} --- อนุมัติแล้ว 59 จาก 61~โครงการ โอนงบร้อยละ~98
        คอขวดอยู่ที่ขั้นตอนการใช้จ่ายในระดับโครงการ
  \item \textbf{ความต่างระหว่างหน่วยงานสูงมาก} --- ตั้งแต่ร้อยละ~96 (มทร.ล้านนา น่าน) ลงมาถึงร้อยละ~0
        บ่งชี้ปัญหาด้านศักยภาพกระบวนการพัสดุรายหน่วยงาน ; ยังมี 7~โครงการที่ยังไม่เบิกจ่ายเลย
        รวม 530,000~บาท
  \item \textbf{ตัวชี้วัดเชิงเครือข่ายและเชิงพาณิชย์ต่ำเชิงโครงสร้าง} --- โครงการส่วนใหญ่ออกแบบ
        เพื่อชุมชนและพื้นที่โครงการหลวง มิได้ออกแบบเพื่อสถานประกอบการหรือการจดทะเบียนทรัพย์สินทางปัญญาตั้งแต่ต้น
\end{enumerate}

\sechead{5. ข้อเสนอต่อผู้บริหาร และเป้าหมาย 55 วันสุดท้าย}

\small
\begin{tabular}{L{5.5cm}L{7.0cm}L{3.4cm}}
\toprule
\rowcolor{headerblue}
\textbf{ข้อเสนอ} & \textbf{การดำเนินการ} & \textbf{หมุดหมาย}\\
\midrule
1. เร่งรัดการเบิกจ่าย & ติดตามโครงการที่เบิกจ่ายต่ำกว่าร้อยละ~30 รายสัปดาห์ พร้อมจัดทีมช่วยจัดทำเอกสารจัดซื้อจัดจ้าง โดยเฉพาะ มทร.ล้านนา ตาก (3~โครงการ 250,000~บาท) และคณะบริหารธุรกิจและศิลปศาสตร์ (1~โครงการ 40,000~บาท) & สิ้น ส.ค. เบิกจ่ายถึงร้อยละ~75\\
\rowcolor{rowblue}
2. กำหนดวินัยการรายงานผล & ให้ผู้รับผิดชอบบันทึกผลรายกิจกรรมพร้อมหลักฐานเข้าระบบติดตามให้ครบ & กลาง ก.ย. ครบ 268~กิจกรรม\\
3. ปิดช่องว่างตัวชี้วัด & เร่งยื่นจดทะเบียนทรัพย์สินทางปัญญาจากผลงานที่พร้อม และลงนามความร่วมมือกับหน่วยงานภายนอกและสถานประกอบการ & ก.ย. 2569\\
\rowcolor{rowblue}
4. ทบทวนการออกแบบตัวชี้วัดปี 2570 & ออกแบบตัวชี้วัดให้สอดคล้องลักษณะงานพัฒนาชุมชนตั้งแต่ต้นปี พร้อมปฏิทินเบิกจ่ายรายไตรมาส & ต้นปีงบประมาณ 2570\\
\bottomrule
\end{tabular}
\normalsize

\vspace{4pt}
\begin{center}
\begin{tikzpicture}
\node[draw=FirstBlue,line width=0.9pt,fill=greybox,rounded corners=3pt,
      inner sep=7pt,text width=16.2cm]{%
  \RaggedRight
  \textbf{\color{FirstBlue}เป้าหมายที่ต้องปิดภายใน 30 กันยายน 2569} \quad
  เบิกจ่ายให้ถึงร้อยละ~90 ของงบจัดสรร $=$ \textbf{7,187,925~บาท} เท่ากับต้องเบิกจ่ายเพิ่มอีก
  \textbf{2,727,593~บาท} ; บันทึกรายงานผลครบ \textbf{268~กิจกรรม} ; ปิดโครงการครบ \textbf{61~โครงการ}};
\end{tikzpicture}
\end{center}

\vspace{2pt}
{\footnotesize\color{black!60}
ตัวเลขทุกจำนวนในเอกสารนี้มาจากฐานข้อมูลระบบติดตามโครงการ ซึ่งปรับปรุงจากไฟล์งบประมาณของกลุ่มแผนงาน
ฉบับวันที่ 4 สิงหาคม 2569 และไฟล์จัดสรรแผนงานและตัวชี้วัด ฉบับวันที่ 20 กรกฎาคม 2569
ตรวจสอบย้อนกลับได้รายโครงการในรายงานความก้าวหน้าฉบับเต็ม (49~หน้า) และที่ระบบติดตามโครงการออนไลน์}

\end{document}
